\section*{Chapitre \ref{ch:g10:measuring-the-earth} --- Mesurer la Terre}

\begin{solution}{exo:g10:measuring-the-earth:1}
Un navire disparaît par la coque d'abord~: sur une mer plate, il ne ferait
que rapetisser, en restant visible tout entier. L'ombre de la Terre sur la
Lune lors d'une éclipse est toujours un arc de cercle~: un disque plat, vu
par la tranche, projetterait une ombre allongée. En allant vers le nord,
l'étoile Polaire monte au-dessus de l'horizon~: sur une Terre plate, sa
hauteur resterait presque la même.
\end{solution}

\begin{solution}{exo:g10:measuring-the-earth:2}
La \omterm{def:g10:measuring-the-earth:geodesy}{géodésie} mesure la forme et les dimensions de la Terre et la position
des points de sa surface, ainsi que leurs variations au cours du temps. La
\omterm{def:g10:measuring-the-earth:triangulation}{triangulation} repère des points par une chaîne de triangles dont tous les
angles sont mesurés~; une seule longueur, la base, est mesurée au sol.
\end{solution}

\begin{solution}{exo:g10:measuring-the-earth:3}
$C = 50 \times 5000 = 250\,000$ stades~; $250\,000 \times \qty{157.2}{m} =
\qty{39300}{km}$.
\end{solution}

\begin{solution}{exo:g10:measuring-the-earth:4}
$111/60 \approx \qty{1.85}{km}$. Quart de méridien~: $90 \times 111 \approx
\qty{10000}{km}$.
\end{solution}

\begin{solution}{exo:g10:measuring-the-earth:5}
$a - b = 6378.137 - 6356.752 = \qty{21.385}{km}$;
$f = 21.385/6378.137 \approx 0.003353 \approx 1/298$.
\end{solution}

\begin{solution}{exo:g10:measuring-the-earth:6}
Au sud~: $\tan\theta_1 = 0.364$, $\theta_1 \approx \qty{20.0}{\degree}$. Au
nord~: $\tan\theta_2 = 0.456$, $\theta_2 \approx \qty{24.5}{\degree}$.
$\alpha = \qty{4.5}{\degree}$, donc $C = (360/4.5) \times 500 = 80 \times 500 =
\qty{40000}{km}$.
\end{solution}

\begin{solution}{exo:g10:measuring-the-earth:7}
$C = 180 - 62 - 71 = \qty{47}{\degree}$.
$AC = 10.00 \sin 71^\circ/\sin 47^\circ = 10.00 \times 0.9455/0.7314 \approx
\qty{12.93}{km}$; $BC = 10.00 \sin 62^\circ/\sin 47^\circ \approx
\qty{12.07}{km}$.
\end{solution}

\begin{solution}{exo:g10:measuring-the-earth:8}
L'écart de longitude est $10 + 20 = \qty{30}{\degree}$. À
\qty{60}{\degree} N, un degré mesure $111.3 \times 0.5 = \qty{55.65}{km}$,
donc la distance vaut $30 \times 55.65 \approx \qty{1670}{km}$. Sur
l'équateur~: $30 \times 111.3 \approx \qty{3340}{km}$, deux fois plus.
\end{solution}

\begin{solution}{exo:g10:measuring-the-earth:9}
La ligne de visée touche la mer en $T$~: une tangente à un cercle est
perpendiculaire au rayon au point de contact, donc l'angle en $T$ est
droit. Pythagore~: $(R+h)^2 = R^2 + d^2$, donc $d^2 = 2Rh + h^2$~; $h$ est
minuscule devant $R$, donc $d \approx \sqrt{2Rh} = \sqrt{2 \times
6\,371\,000 \times 2} \approx \qty{5050}{m}$, environ \qty{5}{km}.
\end{solution}

\begin{solution}{exo:g10:measuring-the-earth:10}
La courbe est la plus pentue aux latitudes moyennes, entre environ
\qty{30}{\degree} et \qty{60}{\degree}, et plate près de l'équateur et du
pôle. Sur une Terre aplatie, la surface est moins courbe près des pôles~;
moins la surface est courbe, plus il faut aller loin pour que la verticale
(et donc la latitude) tourne d'un degré.
\end{solution}

\begin{solution}{exo:g10:measuring-the-earth:11}
$1072.2/9.646 \approx \qty{111.15}{km}$ par degré~; $90 \times 111.15
\approx \qty{10004}{km}$, soit \qty{2}{km} de plus que le vrai quart de
méridien. L'arc est centré vers \qty{46}{\degree}, où un degré est un peu
plus long que le degré moyen de l'équateur au pôle
($10\,002.0/90 \approx \qty{111.13}{km}$).
\end{solution}

\begin{solution}{exo:g10:measuring-the-earth:12}
Avec \qty{7.1}{\degree}~: $C = (360/7.1) \times 5000 \approx 253\,500$ stades
au lieu de $250\,000$~: environ 1.4\,\% de trop (un angle plus petit pour la
même distance signifie une Terre plus grande). La distance, estimée au pas
par des marcheurs le long du Nil, était bien moins précise que l'angle.
\end{solution}

\begin{solution}{exo:g10:measuring-the-earth:13}
Dans $ABC$~: $C = \qty{50}{\degree}$, $AC = 6.000 \sin 75^\circ/\sin 50^\circ
\approx \qty{7.566}{km}$. Dans $ACD$~: l'angle en $D$ vaut $180 - 50 - 68 =
\qty{62}{\degree}$, et $AD = 7.566 \sin 68^\circ/\sin 62^\circ \approx
\qty{7.945}{km}$.
\end{solution}

\begin{solution}{exo:g10:measuring-the-earth:14}
$\cos\varphi = 1/2$ donne $\varphi = \qty{60}{\degree}$~; $\cos\varphi = 0.1$
donne $\varphi \approx \qty{84.3}{\degree}$. Un degré de longitude à la
latitude $\varphi$ est le degré sur l'équateur multiplié par $\cos\varphi$~:
le rapport ne dépend que de l'angle, pas du rayon.
\end{solution}

\begin{solution}{exo:g10:measuring-the-earth:15}
À l'équateur~: $111.195 - 110.574 = \qty{0.621}{km} = \qty{621}{m}$~; au
pôle~: $111.694 - 111.195 = \qty{0.499}{km} = \qty{499}{m}$. Des erreurs de
plusieurs centaines de mètres par degré, qui s'additionnent degré après
degré, dépassent de loin le mètre~: les positions calculées à partir des
latitudes et des longitudes doivent utiliser l'ellipsoïde.
\end{solution}

\begin{solution}{pb:g10:measuring-the-earth:1}
\textbf{1.} $31.20 - 24.09 = \qty{7.11}{\degree}$.

\textbf{2.} $7.11/360 \approx 0.0198 = 1/50.6$, très proche du $1/50$
d'Ératosthène.

\textbf{3.} $C = 788 \times 360/7.11 \approx \qty{39900}{km}$.

\textbf{4.} L'arc de méridien entre les parallèles~: la différence des
angles à midi est la différence de latitude, qui ne mesure que l'écart
nord--sud entre les deux lieux~; la composante est--ouest de la distance ne
joue aucun rôle.

\textbf{5.} Au midi local, le Soleil est dans le plan du méridien, le même
pour les deux lieux~; le même jour, il occupe la même position parmi les
étoiles, de sorte que les deux angles ne diffèrent qu'à cause de la
courbure de la Terre.

\textbf{6.} $180 - 58 - 66 = \qty{56}{\degree}$.

\textbf{7.} $AC = 12.000 \sin 66^\circ/\sin 56^\circ = 12.000 \times
0.9135/0.8290 \approx \qty{13.223}{km}$.

\textbf{8.} $BC = 12.000 \sin 58^\circ/\sin 56^\circ \approx \qty{12.275}{km}$.

\textbf{9.} L'angle en $D$ vaut $180 - 47 - 75 = \qty{58}{\degree}$~;
$AD = 13.223 \sin 75^\circ/\sin 58^\circ \approx \qty{15.061}{km}$.

\textbf{10.} $15.061 \cos 8^\circ \approx \qty{14.914}{km}$.

\textbf{11.} $14.914/111.13 \approx \qty{0.1342}{\degree}$, soit
$0.1342 \times 60 \approx 8.05$ minutes d'arc.

\textbf{12.} $111.51 - 110.57 = \qty{0.94}{km}$: \qty{940}{m}.

\textbf{13.} Près du pôle, la verticale tourne plus lentement avec la
distance~: la surface y est moins courbe. La Terre est aplatie aux pôles et
renflée à l'équateur.

\textbf{14.} $a - b = \qty{21.385}{km}$~; $f = 21.385/6378.137 \approx 1/298$.

\textbf{15.} Le fil à plomb est perpendiculaire à la surface de
l'ellipsoïde. Sur une ellipse, la perpendiculaire à la courbe ne passe par
le centre qu'à l'équateur et aux pôles~; à \qty{45}{\degree}, elle coupe
l'axe sous le centre.

\textbf{16.} $10\,000\,000$ mètres, soit \qty{10000}{km}.

\textbf{17.} $10\,001\,966/10\,000\,000 \approx \qty{1.0001966}{m}$.

\textbf{18.} Trop court~: la vraie dix-millionième partie vaut $1.0001966$
fois le mètre de 1799, donc le mètre de 1799 est trop court d'environ
$1.97 \times 10^{-4}$ de sa longueur, soit 0.02\,\%.

\textbf{19.} Trop court. La barre est la dix-millionième partie du quart de
méridien tel qu'on le calculait dans les années 1790~; comme la barre est
plus courte que la dix-millionième partie du vrai quart de méridien, le
quart de méridien calculé était plus court que le vrai.

\textbf{20.} $1.0001966 - 1 = 0.0001966$ m~: le premier mètre était trop
court d'environ \textbf{\qty{0.2}{mm}} --- à peu près deux feuilles de
papier sur un mètre, une précision étonnante pour une mesure faite avec des
lunettes et des angles à travers tout un pays.
\end{solution}
